\documentclass[aspectratio=169,11pt]{beamer}

\usepackage[utf8]{inputenc}
\usepackage[T5]{fontenc}
\usepackage[vietnamese,english]{babel}
\usepackage{lmodern}
\usepackage{graphicx}
\usepackage{booktabs}
\usepackage{array}
\usepackage{hyperref}

\definecolor{TimelyBlue}{HTML}{173B57}
\definecolor{TimelyTeal}{HTML}{168C8C}
\definecolor{TimelyLight}{HTML}{EEF4F5}
\definecolor{TimelyGray}{HTML}{4B5563}

\setbeamercolor{normal text}{fg=TimelyGray,bg=white}
\setbeamercolor{structure}{fg=TimelyTeal}
\setbeamercolor{frametitle}{fg=white,bg=TimelyBlue}
\setbeamercolor{title}{fg=TimelyBlue}
\setbeamercolor{subtitle}{fg=TimelyGray}
\setbeamercolor{block title}{fg=white,bg=TimelyTeal}
\setbeamercolor{block body}{fg=TimelyGray,bg=TimelyLight}
\setbeamercolor{alerted text}{fg=TimelyTeal}

\setbeamertemplate{navigation symbols}{}
\setbeamertemplate{itemize item}{\color{TimelyTeal}\small$\blacktriangleright$}
\setbeamertemplate{itemize subitem}{\color{TimelyTeal}\scriptsize$\bullet$}
\setbeamertemplate{footline}{%
  \leavevmode%
  \hbox{%
    \begin{beamercolorbox}[wd=.86\paperwidth,ht=2.7ex,dp=1.2ex,leftskip=1em]{author in head/foot}%
      \color{TimelyGray}\scriptsize TimelyMT \hspace{0.5em}|\hspace{0.5em} Policy quyết định thời điểm dịch
    \end{beamercolorbox}%
    \begin{beamercolorbox}[wd=.14\paperwidth,ht=2.7ex,dp=1.2ex,center]{date in head/foot}%
      \color{TimelyBlue}\scriptsize\insertframenumber/\inserttotalframenumber
    \end{beamercolorbox}%
  }%
}
\setbeamertemplate{headline}{}
\setbeamersize{text margin left=0.7cm,text margin right=0.7cm}

\newcommand{\slidefigure}[2][]{%
  \IfFileExists{figures/#2}{%
    \includegraphics[#1]{figures/#2}%
  }{%
    \includegraphics[#1]{../figures/#2}%
  }%
}
\newcommand{\contextreal}{\texttt{REAL\_CONTEXT}}
\newcommand{\contextzero}{\texttt{ZERO\_CONTEXT}}

\title{Nghiên cứu mô hình quyết định thời điểm dịch\\
cho hệ thống dịch cabin tự động Anh--Việt\\
với độ trễ thấp}
\subtitle{Môn: Xử lý ngôn ngữ tự nhiên}
\author{Sinh viên: Đặng Quang Minh \quad | \quad MSSV: 24022397}
\date{\url{https://github.com/MinhCYB/TimelyMT}}

\begin{document}

% 1. Title
\begin{frame}[plain]
  \titlepage
  \vspace{-0.8em}
  \begin{center}
    \color{TimelyTeal}\rule{0.62\textwidth}{1.2pt}
  \end{center}
\end{frame}

% 2. Problem
\begin{frame}{Bài toán: khi nào nên phát hành bản dịch?}
  \begin{columns}[T,onlytextwidth]
    \begin{column}{0.59\textwidth}
      \begin{itemize}
        \setlength\itemsep{0.8em}
        \item Dịch theo luồng chỉ quan sát từng phần của câu nguồn.
        \item Policy phải quyết định \textbf{chờ thêm} hay \textbf{COMMIT} ngay.
        \item Một khi đã \textbf{COMMIT}, phần dịch đó không thể sửa lại.
      \end{itemize}
    \end{column}
    \begin{column}{0.37\textwidth}
      \begin{block}{Mục tiêu dài hạn}
        Dịch cabin tự động Anh--Việt với độ trễ thấp.
      \end{block}
      \vspace{0.4em}
       \small\textbf{Phạm vi hiện tại:} policy trên luồng token và đồng hồ mô phỏng; chưa phải hệ thống speech-to-speech end-to-end.
    \end{column}
  \end{columns}
\end{frame}

% 3. Pipeline
\begin{frame}{Luồng xử lý của TimelyMT}
  \begin{center}
    \slidefigure[width=0.97\textwidth,height=0.75\textheight,keepaspectratio]{figure-1-streaming-pipeline.png}
  \end{center}
  \vspace{-0.7em}
  \centering EnViT5 tạo bản dịch ứng viên; \alert{policy quyết định khi nào bản dịch đó nên được dịch}.
\end{frame}

% 4. Dataset and experimental protocol
\begin{frame}{Dataset và thiết kế thực nghiệm}
  \begin{columns}[T,onlytextwidth]
    \begin{column}{0.235\textwidth}
      \begin{beamercolorbox}[rounded=true,sep=0.45em,center]{block body}
        {\LARGE\bfseries\color{TimelyBlue}17}\par
        \scriptsize bài TED
      \end{beamercolorbox}
    \end{column}
    \begin{column}{0.235\textwidth}
      \begin{beamercolorbox}[rounded=true,sep=0.45em,center]{block body}
        {\LARGE\bfseries\color{TimelyBlue}41.7K}\par
        \scriptsize token nguồn tiếng Anh
      \end{beamercolorbox}
    \end{column}
    \begin{column}{0.235\textwidth}
      \begin{beamercolorbox}[rounded=true,sep=0.45em,center]{block body}
        {\LARGE\bfseries\color{TimelyBlue}278 phút}\par
        \scriptsize 
      \end{beamercolorbox}
    \end{column}
    \begin{column}{0.235\textwidth}
      \begin{beamercolorbox}[rounded=true,sep=0.45em,center]{block body}
        {\LARGE\bfseries\color{TimelyBlue}12 / 3 / 2}\par
        \scriptsize TRAIN / DEV / TEST
      \end{beamercolorbox}
    \end{column}
  \end{columns}

  \vspace{0.35em}
  \begin{columns}[T,onlytextwidth]
    \begin{column}{0.485\textwidth}
      \begin{block}{PHÂN CHIA DỮ LIỆU}
        \centering
        \begin{tabular}{@{}>{\centering\arraybackslash}p{0.29\linewidth}>{\centering\arraybackslash}p{0.29\linewidth}>{\centering\arraybackslash}p{0.29\linewidth}@{}}
          \textbf{TRAIN} & \textbf{DEV} & \textbf{TEST} \\
          {\Large\color{TimelyBlue}\textbf{12}} & {\Large\color{TimelyBlue}\textbf{3}} & {\Large\color{TimelyBlue}\textbf{2}} \\
          \scriptsize bài & \scriptsize bài & \scriptsize bài \\
          \scriptsize huấn luyện policy & \scriptsize phân tích / lựa chọn & \scriptsize đánh giá sau khi cố định thiết kế
        \end{tabular}

        \vspace{0.4em}
        \small\bfseries Chia theo toàn bài \;\textperiodcentered\; Theo người nói \;\textperiodcentered\; Không trùng người nói
      \end{block}
    \end{column}
    \begin{column}{0.485\textwidth}
      \begin{block}{MÔ PHỎNG STREAMING}
        \centering\small
        Transcript tiếng Anh từ TED

        \vspace{-0.15em}
        {\color{TimelyTeal}$\downarrow$}\par
        Luồng token từ vựng

        \vspace{-0.15em}
        {\color{TimelyTeal}$\downarrow$}\par
        \textbf{đồng hồ nguồn mô phỏng} \;\alert{(2,5 token / giây)}

        \vspace{-0.15em}
        {\color{TimelyTeal}$\downarrow$}\par
        \textbf{WAIT / LISTEN / COMMIT}
      \end{block}
    \end{column}
  \end{columns}

  \vspace{0.15em}
  \centering\scriptsize
  \textbf{Bản tham chiếu tiếng Việt} $\rightarrow$ chỉ dùng để căn chỉnh và đánh giá\par
  \vspace{0.15em}
  22.018 trạng thái TRAIN $\neq$ 22.018 bài nói độc lập
\end{frame}

% 5. Research journey
\begin{frame}{Hành trình nghiên cứu: V1 $\rightarrow$ V2 $\rightarrow$ P3\_GLOBAL}
  \vspace{0.1em}
  \begin{columns}[T,onlytextwidth]
    \begin{column}{0.30\textwidth}
      \begin{beamercolorbox}[rounded=true,sep=0.45em,center]{block body}
        {\Large\bfseries\color{TimelyBlue}V1}\par
        \vspace{0.15em}\textbf{\small Thiết lập baseline nhân quả}\par
        \vspace{0.3em}{\scriptsize\color{TimelyGray}TF-IDF + 11 đặc trưng số}\par
        \vspace{0.3em}\rule{0.8\linewidth}{0.4pt}\par
        \vspace{0.3em}{\bfseries\color{TimelyTeal}Kết quả}\par
        \vspace{0.15em}\textbf{\footnotesize P1 @ 0.60}\;{\scriptsize được chọn trên DEV}\par
        \vspace{0.3em}{\scriptsize Xác lập pipeline causal\par WAIT / LISTEN / COMMIT}
      \end{beamercolorbox}
    \end{column}
    \begin{column}{0.04\textwidth}
      \vspace{4.2em}\centering\color{TimelyTeal}$\rightarrow$
    \end{column}
    \begin{column}{0.30\textwidth}
      \begin{beamercolorbox}[rounded=true,sep=0.45em,center]{block body}
        {\Large\bfseries\color{TimelyBlue}V2}\par
        \vspace{0.15em}\textbf{\small Mở rộng biểu diễn}\par
        \vspace{0.3em}{\scriptsize\color{TimelyGray}MiniLM + lịch sử nguồn / đích đã quan sát}\par
        \vspace{0.3em}\rule{0.8\linewidth}{0.4pt}\par
        \vspace{0.3em}{\bfseries\color{TimelyTeal}Kết quả}\par
        \vspace{0.15em}\textbf{\footnotesize P2 @ 0.50}\;{\scriptsize được chọn trong V2}\par
        \vspace{0.3em}{\scriptsize\bfseries\color{TimelyTeal}DEV không xác lập\par ưu thế chất lượng so với V1}
      \end{beamercolorbox}
    \end{column}
    \begin{column}{0.04\textwidth}
      \vspace{4.2em}\centering\color{TimelyTeal}$\rightarrow$
    \end{column}
    \begin{column}{0.30\textwidth}
      \begin{beamercolorbox}[rounded=true,sep=0.45em,center]{block body}
        {\Large\bfseries\color{TimelyBlue}P3\_GLOBAL}\par
        \vspace{0.15em}\textbf{\small Kiểm tra prepared context}\par
        \vspace{0.3em}{\scriptsize\color{TimelyGray}+ 384 chiều prepared context}\par
        \vspace{0.3em}\rule{0.8\linewidth}{0.4pt}\par
        \vspace{0.3em}{\bfseries\color{TimelyTeal}Kết quả}\par
        \vspace{0.15em}\textbf{\footnotesize REAL / ZERO}\;{\scriptsize làm thay đổi COMMIT}\par
        \vspace{0.3em}{\scriptsize Lợi ích chất lượng / độ trễ\par chưa ổn định}
      \end{beamercolorbox}
    \end{column}
  \end{columns}
  \vfill
  \centering\footnotesize\color{TimelyGray}
  Mỗi giai đoạn mở rộng một câu hỏi nghiên cứu; phiên bản sau không mặc định tốt hơn phiên bản trước.
\end{frame}

% 6. V1 causal baseline
\begin{frame}{V1 --- baseline nhân quả}
  \begin{columns}[T,onlytextwidth]
    \begin{column}{0.48\textwidth}
      \begin{block}{Biểu diễn}
        \scriptsize
        Biểu diễn văn bản \textbf{TF--IDF} unigram/bigram\par
        \vspace{0.2em}
        \textbf{11 đặc trưng số} theo nguyên tắc nhân quả\par
        \vspace{0.2em}
        Logistic regression có cân bằng lớp
      \end{block}
      \vspace{0.25em}
      \scriptsize
      \begin{tabular}{@{}>{\bfseries}p{0.10\linewidth}p{0.84\linewidth}@{}}
        P0 & nguồn hiện tại \\[0.25em]
        P1 & nguồn hiện tại\\
           & + nguồn đã COMMIT trước đó \\[0.25em]
        P2 & nguồn hiện tại\\
           & + nguồn đã COMMIT trước đó\\
           & + đích đã sinh và COMMIT trước đó
      \end{tabular}
      \vspace{0.5em}
      \begin{beamercolorbox}[rounded=true,sep=0.45em,center]{block body}
        Cấu hình V1 được chọn: \textbf{\texttt{learned\_P1\_0.60}}\par
        \scriptsize được chọn trên DEV; không phải khẳng định tối ưu toàn cục
      \end{beamercolorbox}
    \end{column}
    \begin{column}{0.49\textwidth}
      \begin{block}{Quy tắc khi chạy: WAIT / LISTEN / COMMIT}
        \scriptsize
        \begin{tabular}{@{}>{\bfseries}p{0.32\linewidth}p{0.61\linewidth}@{}}
          $<4$ token nguồn & WAIT; không gọi EnViT5; không tính xác suất policy \\[0.45em]
          $\geq4$ token nguồn & EnViT5 tạo ứng viên + policy ra quyết định \\[0.45em]
          48 token nguồn & COMMIT cưỡng bức \\[0.45em]
          cuối bài nói & phần nguồn còn lại được COMMIT.
        \end{tabular}
      \end{block}
      \vspace{0.55em}
      \begin{beamercolorbox}[rounded=true,sep=0.55em,center]{block body}
        \textbf{Phần đã COMMIT không thể sửa lại.}
      \end{beamercolorbox}
    \end{column}
  \end{columns}
\end{frame}

% 7. V1 supervision
\begin{frame}{V1 --- tạo dữ liệu giám sát LISTEN / COMMIT}
  \begin{columns}[T,onlytextwidth]
    \begin{column}{0.58\textwidth}
      \begin{block}{TẠO NHÃN NGOẠI TUYẾN}
        \centering\scriptsize
        bản dịch ứng viên hiện tại $h_t$\par
        \vspace{-0.15em}{\color{TimelyTeal}$\downarrow$}\par
        ứng viên sau khi thêm 1 token nguồn\par
        \vspace{-0.15em}{\color{TimelyTeal}$\downarrow$}\par
        ứng viên sau khi thêm 2 token nguồn\par
        \vspace{-0.15em}{\color{TimelyTeal}$\downarrow$}\par
        độ ổn định khi nhìn thêm tương lai\par
        \vspace{-0.15em}{\color{TimelyTeal}$\downarrow$}\par
        \textbf{pseudo-label LISTEN / COMMIT}
      \end{block}
      \vspace{0.25em}
      \centering\scriptsize
      $\displaystyle \mathit{future\_stability\_ratio}
        = \frac{\operatorname{LCP}(h_t,h_{t+1},h_{t+2})}{\operatorname{len}(h_t)}$

      \vspace{0.35em}
      $\text{độ ổn định}\geq \text{threshold} \Rightarrow \textbf{COMMIT}$
      \quad ngược lại $\Rightarrow \textbf{LISTEN}$
    \end{column}
    \begin{column}{0.39\textwidth}
      \begin{alertblock}{Ranh giới nhân quả}
        \scriptsize
        \textbf{TẠO NHÃN NGOẠI TUYẾN}\par
        có thể dùng ứng viên tương lai $+1$ / $+2$

        \vspace{0.55em}
        \textbf{KHI CHẠY TRỰC TUYẾN}\par
        chỉ dùng thông tin đã quan sát
      \end{alertblock}
      \scriptsize
      \begin{itemize}
        \setlength\itemsep{0.25em}
        \item Nhìn trước tối đa: \textbf{2 lần cập nhật token nguồn}.
        \item Thông tin tương lai chỉ dùng để tạo nhãn giám sát TRAIN / DEV ngoại tuyến.
        \item Thông tin tương lai không bao giờ đi vào đầu vào policy khi chạy.
        \item \textbf{Không tạo pseudo-label cho TEST.}
      \end{itemize}
      \vspace{0.15em}
      \begin{beamercolorbox}[rounded=true,sep=0.45em,center]{block body}
        \textbf{22.018 trạng thái policy TRAIN}\par
        \vspace{0.15em}
        \textbf{19.075 LISTEN}\quad |\quad \textbf{2.943 COMMIT}\par
        \vspace{0.2em}
        \scriptsize Tín hiệu oracle chỉ dùng khi huấn luyện
      \end{beamercolorbox}
    \end{column}
  \end{columns}
\end{frame}

% 8. V2 semantic policy
\begin{frame}{V2 --- policy ngữ nghĩa}
  \small
  V1 dùng biểu diễn từ vựng thưa. V2 thử nghiệm biểu diễn ngữ nghĩa phong phú hơn, nhưng vẫn giữ \textbf{ranh giới nhân quả khi chạy} và \textbf{cùng dữ liệu giám sát của V1}.
  \vspace{0.35em}
  \begin{columns}[T,onlytextwidth]
    \begin{column}{0.54\textwidth}
      \begin{block}{Các biến thể mở rộng dần}
        \scriptsize
        \begin{tabular}{@{}>{\bfseries}p{0.07\linewidth}p{0.58\linewidth}>{\raggedleft\arraybackslash}p{0.24\linewidth}@{}}
          P0 & embedding nguồn hiện tại & $384+11=395$ \\
          \addlinespace[0.25em]
          P1 & \shortstack[l]{nguồn hiện tại +\\nguồn đã COMMIT trước đó} & $2\times384+11=779$ \\
          \addlinespace[0.25em]
          P2 & \shortstack[l]{nguồn hiện tại +\\nguồn đã COMMIT trước đó +\\đích đã sinh trước đó} & $3\times384+11=1163$
        \end{tabular}
      \end{block}
      \scriptsize
      \textbf{Bộ mã hóa:} \texttt{sentence-transformers/}\par
      \texttt{paraphrase-multilingual-MiniLM-L12-v2}\par
      được cố định, embedding 384 chiều
    \end{column}
    \begin{column}{0.43\textwidth}
      \begin{beamercolorbox}[rounded=true,sep=0.55em,center]{block body}
        \textbf{Policy:} MLP phi tuyến\par
        \vspace{0.4em}
        Cấu hình V2 được chọn:\par
        \textbf{\texttt{v2\_P2\_0.50}}
      \end{beamercolorbox}
      \vspace{0.45em}
      \begin{alertblock}{Kết quả trên DEV}
        \centering\bfseries DEV KHÔNG xác lập ưu thế về chất lượng của V2 so với V1.
      \end{alertblock}
    \end{column}
  \end{columns}
  \vfill
  \centering\scriptsize
  V2 cho thấy có thể dùng biểu diễn ngữ nghĩa và lịch sử nguồn/đích đã quan sát\par
  mà vẫn giữ ranh giới nhân quả, đồng thời tạo nền biểu diễn để mở rộng sang P3.
\end{frame}

% 9. P3_GLOBAL
\begin{frame}{P3\_GLOBAL --- prepared context trước bài nói}
  \begin{columns}[T,onlytextwidth]
    \begin{column}{0.55\textwidth}
      \begin{block}{Đầu vào policy: 1547 chiều}
        \scriptsize
        \begin{tabular}{@{}rl@{}}
          384 & nguồn hiện tại \\
          + 384 & nguồn đã COMMIT trước đó \\
          + 384 & đích đã sinh trước đó \\
          + 384 & prepared context \\
          + 11 & đặc trưng số \\
          \midrule
          1547 & tổng
        \end{tabular}
      \end{block}
      \scriptsize
      \textbf{Prepared context chỉ được đưa vào policy.}\par
      EnViT5 vẫn được cố định, chỉ nhận nguồn và không nhận prepared context.
      \vspace{0.45em}
      \begin{beamercolorbox}[rounded=true,sep=0.45em,center]{block body}
        \textbf{MLP:} $1547 \rightarrow 256 \rightarrow 64 \rightarrow 1$\par
        \scriptsize GELU; dropout 0.20 / 0.10; đầu ra $p(\mathrm{COMMIT})$
      \end{beamercolorbox}
    \end{column}
    \begin{column}{0.42\textwidth}
      \begin{alertblock}{Huấn luyện từ đầu}
        \scriptsize P3 không khởi tạo từ V2 P2. P3 tái sử dụng:
        \begin{itemize}
          \item cùng 22.018 dòng dữ liệu giám sát TRAIN của V1;
          \item cùng 11 định nghĩa đặc trưng số;
          \item cùng quy trình EnViT5 cố định để tạo bản dịch ứng viên.
        \end{itemize}
      \end{alertblock}
      \centering\small
      \textbf{P3 $-$ P2 chỉ mang tính mô tả, không phải ablation sạch để cô lập ảnh hưởng của prepared context.}
    \end{column}
  \end{columns}
  \vfill
  \begin{beamercolorbox}[rounded=true,sep=0.5em,center]{block title}
    \textbf{Thông tin có trước bài nói có thể làm thay đổi thời điểm COMMIT không?}
  \end{beamercolorbox}
\end{frame}

% 10. Prepared-context provenance
\begin{frame}{Prepared context: nguồn gốc và kiểm soát rò rỉ}
  \begin{columns}[T,onlytextwidth]
    \begin{column}{0.34\textwidth}
      \begin{block}{Có sẵn trước bài nói}
        \small
        \begin{itemize}
          \setlength\itemsep{0.35em}
          \item Bài báo
          \item Trang dự án
          \item Tài liệu nền
        \end{itemize}
      \end{block}
      \scriptsize\textbf{Không dùng} transcript hoặc bản dịch tham chiếu.
    \end{column}
    \begin{column}{0.35\textwidth}
      \begin{block}{Điều kiện hợp lệ}
        \scriptsize\ttfamily
        SAFE\_PRETALK\_CONFIRMED\par
        available\_before\_talk = true\par
        transcript\_used = false\par
        reference\_used = false
      \end{block}
    \end{column}
    \begin{column}{0.27\textwidth}
      \begin{beamercolorbox}[rounded=true,sep=0.45em]{block body}
        \scriptsize
        \textbf{Phạm vi dữ liệu}\par
        12 pool TRAIN\par
        3 pool DEV\par
        8 nguồn hợp lệ\par
        5 bài TRAIN có context khác 0\par
        1 bài DEV có context khác 0
      \end{beamercolorbox}
    \end{column}
  \end{columns}
  \vspace{0.55em}
  \begin{beamercolorbox}[rounded=true,sep=0.55em,center]{block body}
    \textbf{DEV:} Sims = context thực \quad | \quad Jeff / Luis = đối chứng context rỗng
  \end{beamercolorbox}
  \vfill
  \centering\alert{Prepared context chỉ đi vào policy, không bao giờ đi vào EnViT5.}
\end{frame}

% 11. Why controlled ablation
\begin{frame}{Vì sao cần ablation có kiểm soát?}
  \begin{columns}[c,onlytextwidth]
    \begin{column}{0.44\textwidth}
      \begin{block}{So sánh trực tiếp P3--P2}
        \small P3 và P2 được huấn luyện riêng. Chênh lệch có thể đến từ quá trình huấn luyện, kiến trúc và prepared context.
        \vspace{0.45em}

        Vì vậy, \textbf{P3--P2 không cô lập riêng ảnh hưởng của prepared context}.
      \end{block}
    \end{column}
    \begin{column}{0.08\textwidth}
      \centering\Large\color{TimelyTeal}$\Longrightarrow$
    \end{column}
    \begin{column}{0.44\textwidth}
      \begin{block}{Câu hỏi nhân quả hẹp hơn}
        \small Với \textbf{cùng một P3 checkpoint} và cùng rollout, quyết định COMMIT có đổi khi chỉ lát prepared context bị loại bỏ?
      \end{block}
    \end{column}
  \end{columns}
  \vfill
  \centering\small Đây là lý do cần thí nghiệm \contextreal{} / \contextzero{}.
\end{frame}

% 12. Controlled ablation
\begin{frame}{\contextreal{} vs \contextzero{}}
  \begin{columns}[T,onlytextwidth]
    \begin{column}{0.68\textwidth}
      \centering
      \slidefigure[width=\textwidth,height=0.73\textheight,keepaspectratio]{figure-3-controlled-ablation.png}
    \end{column}
    \begin{column}{0.29\textwidth}
      \begin{block}{Chỉ thay đổi}
        \centering\small lát prepared context

        \vspace{0.2em}
        \Large\texttt{[1152:1536]}
      \end{block}
      \vspace{0.45em}
      \scriptsize
      \textbf{Giữ nguyên:}
      \begin{itemize}
        \item P3 checkpoint
        \item luồng nguồn
        \item ngưỡng
        \item EnViT5
      \end{itemize}
    \end{column}
  \end{columns}
\end{frame}

% 13. Evaluation metrics
\begin{frame}{Các chỉ số đánh giá: chất lượng và độ trễ}
  \begin{columns}[T,onlytextwidth]
    \begin{column}{0.485\textwidth}
      \begin{block}{CHẤT LƯỢNG DỊCH}
        \scriptsize
        \textbf{\normalsize BLEU $\uparrow$}\par
        \vspace{0.1em}
        Độ trùng khớp $n$-gram giữa bản dịch hệ thống và reference (càng cao càng tốt).\par
        \vspace{0.15em}
        \textit{Lưu ý:} Có brevity penalty để hạn chế bản dịch quá ngắn.

        \vspace{0.45em}
        \textbf{\normalsize chrF2 $\uparrow$}\par
        \vspace{0.1em}
        F-score trên character $n$-gram; $\beta = 2$ nên nhấn mạnh recall (càng cao càng tốt).\par
        \vspace{0.15em}
        \textit{Đặc điểm:} Nhạy với mức độ giống nhau ở cấp ký tự, ít phụ thuộc tách từ hơn BLEU.
      \end{block}
    \end{column}
    \begin{column}{0.485\textwidth}
      \begin{block}{ĐỘ TRỄ / HÀNH VI POLICY}
        \scriptsize
        \textbf{\normalsize AL $\downarrow$ \textnormal{(Average Lagging)}}\par
        \vspace{0.1em}
        Đo policy trung bình đi sau source bao nhiêu theo vị trí source token (càng thấp càng ít trễ).\par
        \vspace{0.15em}
        \textit{Lưu ý:} AL âm $\neq$ ``độ trễ thời gian âm''; đây là giá trị tương đối so với ideal policy.

        \vspace{0.35em}
        \textbf{\normalsize LAAL $\downarrow$ \textnormal{(Length-Adaptive AL)}}\par
        \vspace{0.1em}
        Biến thể của AL có hiệu chỉnh theo độ dài output, giảm khả năng hệ thống được ``thưởng'' do khác biệt độ dài.

        \vspace{0.35em}
        \textbf{\normalsize COMMIT}\par
        \vspace{0.1em}
        Số đơn vị dịch được policy phát hành không thể sửa lại. Phản ánh cách policy phân đoạn và phát hành bản dịch; nhiều hay ít không tự động tốt hơn.
      \end{block}
    \end{column}
  \end{columns}
  \vfill
  \centering\scriptsize\color{TimelyGray}
  \textbf{Lưu ý:} Latency trong nghiên cứu này được tính trên vị trí lexical source token / simulated source clock, không phải acoustic latency thực.
\end{frame}

% 14. DEV result
\begin{frame}{Kết quả DEV: prepared context thay đổi COMMIT, lợi ích chưa ổn định}
  \begin{columns}[T,onlytextwidth]
    \begin{column}{0.47\textwidth}
      \begin{block}{Trên lưới năm ngưỡng}
        REAL và ZERO tạo số lượng \textbf{COMMIT} khác nhau ở cả năm ngưỡng.
      \end{block}
      \small
      Đây là bằng chứng nhân quả rằng prepared context có thể thay đổi hành vi COMMIT của P3 đã huấn luyện.
      \vspace{0.35em}

      Tuy nhiên, BLEU, chrF2 và các chỉ số độ trễ không cải thiện nhất quán.
    \end{column}
    \begin{column}{0.49\textwidth}
      \centering
      \textbf{Sims 0.60: một trường hợp mô tả}

      \vspace{0.35em}
      \begin{tabular}{>{\raggedright\arraybackslash}p{0.45\linewidth}r}
        \toprule
        Chỉ số & Giá trị \\
        \midrule
        $\Delta$ BLEU & $+0.95$ \\
        $\Delta$ chrF2 & $+0.63$ \\
        $\Delta$ COMMIT & $+19$ \\
        \bottomrule
      \end{tabular}

      \vspace{0.25em}
      \scriptsize Số COMMIT REAL/ZERO: 318 vs 299\par
      \alert{Mô tả, không phải ngưỡng tối ưu.}
    \end{column}
  \end{columns}
  \vfill
  \begin{alertblock}{Giới hạn kết luận}
    \texttt{prepared-global-v0} chưa chứng minh cải thiện ổn định về chất lượng hoặc độ trễ.
  \end{alertblock}
\end{frame}

% 15. Held-out TEST
\begin{frame}{TEST: kiểm tra sau khi cố định thiết kế}
  \begin{beamercolorbox}[rounded=true,sep=0.25em,center]{block body}
    \scriptsize
    \textbf{2 TEST talks} \;\textperiodcentered\; Mở đúng một lần sau freeze \;\textperiodcentered\; Không tuning sau TEST \;\textperiodcentered\; Không chọn lại model / threshold\par
    \tiny\color{TimelyGray}P3 TEST = \contextzero{} vì không có approved prepared-context pool cho TEST.
  \end{beamercolorbox}
  \vspace{0.15em}

  \begin{columns}[T,onlytextwidth]
    \begin{column}{0.485\textwidth}
      \begin{block}{So sánh V1 / V2 được chọn}
        \centering\scriptsize
        \begin{tabular}{lcccc}
          \toprule
          Cấu hình & BLEU & chrF2 & AL & LAAL \\
          \midrule
          V1 P1 @ 0.60 & 23.38 & 56.29 & 1.22 & 41.26 \\
          V2 P2 @ 0.50 & 21.05 & 55.34 & -0.07 & 40.15 \\
          \bottomrule
        \end{tabular}

        \vspace{0.2em}
        {\bfseries\color{TimelyTeal}V1 $-$ V2 trên TEST: $+2.32$ BLEU \quad $+0.96$ chrF2}
      \end{block}
      \vspace{0.1em}
      \scriptsize
      \begin{itemize}
        \setlength\itemsep{0.15em}
        \item V1 selected $>$ V2 selected về chất lượng trên TEST ($+2.32$ BLEU / $+0.96$ chrF2).
        \item Chất lượng TEST của các cấu hình thấp hơn các giá trị DEV tương ứng.
      \end{itemize}
    \end{column}

    \begin{column}{0.485\textwidth}
      \begin{block}{Lưới ngưỡng P3 \contextzero{}}
        \centering\scriptsize
        \begin{tabular}{ccccc}
          \toprule
          Ngưỡng & BLEU & chrF2 & AL & LAAL \\
          \midrule
          0.30 & 19.26 & 55.38 & 10.55 & 13.59 \\
          0.40 & 19.25 & 55.47 & 7.36 & 16.98 \\
          0.50 & 21.98 & 56.71 & 13.69 & 20.38 \\
          0.60 & 20.57 & 55.11 & 18.01 & 55.35 \\
          0.70 & 22.79 & 55.61 & 25.09 & 56.61 \\
          \bottomrule
        \end{tabular}
      \end{block}
      \vspace{0.1em}
      \scriptsize
      \begin{itemize}
        \setlength\itemsep{0.15em}
        \item P3 vẫn thể hiện trade-off theo threshold, nhưng TEST không đánh giá lợi ích prepared context vì P3 chỉ chạy \contextzero{}.
      \end{itemize}
    \end{column}
  \end{columns}

  \vfill
  \centering\scriptsize\color{TimelyGray}
\end{frame}

% 16. Divergence cascade
\begin{frame}{Event 131--132: từ một COMMIT đến phân kỳ về sau}
  \begin{columns}[T,onlytextwidth]
    \begin{column}{0.57\textwidth}
      \centering
      \slidefigure[width=\textwidth,height=0.69\textheight,keepaspectratio]{figure-4-divergence-cascade.png}
    \end{column}
    \begin{column}{0.40\textwidth}
      \scriptsize
      \textbf{Event 131}
      \begin{itemize}
        \item Cùng đoạn nguồn, cùng bản dịch ứng viên.
        \item REAL: $p = 0.631034$ $\rightarrow$ COMMIT.
        \item ZERO: $p = 0.575424$ $\rightarrow$ LISTEN.
      \end{itemize}
      \vspace{0.25em}
      \textbf{Event 132}
      \begin{itemize}
        \item Hai phiên bắt đầu có bộ đệm khác nhau.
        \item Từ đây, EnViT5 có thể nhận các đoạn nguồn khác nhau.
      \end{itemize}
      \vspace{0.25em}
      \alert{Prepared context làm thay đổi thời điểm COMMIT; đây không phải bằng chứng về chất lượng dịch tốt hơn.}
    \end{column}
  \end{columns}
\end{frame}

% 17. Demo and conclusion
\begin{frame}{Demo và kết luận}
  \begin{columns}[T,onlytextwidth]
    \begin{column}{0.48\textwidth}
      \centering
      \slidefigure[width=\textwidth,height=0.55\textheight,keepaspectratio]{figure-5-interactive-demo.png}

      \vspace{0.3em}
      \scriptsize\textbf{Trace replay từ causal streaming rollout}\par
      \vspace{0.15em}
      \tiny\color{TimelyGray}Demo phát lại trace của một causal streaming rollout để quan sát quá trình phân kỳ giữa REAL và ZERO.
    \end{column}
    \begin{column}{0.49\textwidth}
      \footnotesize
      \begin{itemize}
        \setlength\itemsep{0.25em}
        \item \textbf{V1:} Thiết lập baseline nhân quả, xác lập ranh giới và ngữ nghĩa WAIT / LISTEN / COMMIT.
        \item \textbf{V2:} Mở rộng biểu diễn ngữ nghĩa và lịch sử đã quan sát; DEV không xác lập ưu thế chất lượng so với V1.
        \item \textbf{P3:} Đặt câu hỏi nghiên cứu về prepared context; REAL/ZERO làm thay đổi COMMIT nhưng lợi ích chất lượng / độ trễ chưa ổn định.
        \item \textbf{Hệ thống:} Mô hình nghiên cứu trên luồng token, chưa phải hệ thống speech-to-speech end-to-end.
        \item \textbf{TEST:} Mở một lần duy nhất sau khi cố định thiết kế; không điều chỉnh lại theo TEST.
      \end{itemize}
    \end{column}
  \end{columns}
\end{frame}

\end{document}
